% ============================================================================
%  The iliac crest as a level landmark — clinical measurements paper.
%
%  TARGET: Spine Surgery and Related Research (SSRR).
%    - No article processing charge.
%    - PubMed, Scopus, DOAJ, Embase indexed.
%    - ~15 weeks submission to publication.
%    - Official journal of the Japanese Society for Spine Surgery and Related Research;
%      scope covers clinical spine research including imaging morphometry.
%
%  WHY THIS IS A SEPARATE PAPER FROM THE DATASET ARTICLE. The Medical Physics dataset
%  article may not contain hypothesis testing or generalisable conclusions -- that is
%  policy, not preference. This paper is where the clinical claim lives, and the claim is
%  a single sentence: the landmark surgeons use to find L4-5 fails hardest in exactly the
%  patients whose levels are already ambiguous.
%
%  STATUS: complete draft. Needs a literature introduction expanded by a co-author, and
%  someone to argue with the interpretation in the discussion.
% ============================================================================
\documentclass[11pt,a4paper]{article}

\usepackage[margin=1in]{geometry}
\usepackage{graphicx}
\usepackage{booktabs}
\usepackage{siunitx}
\usepackage{authblk}
\usepackage[colorlinks=true,allcolors=blue]{hyperref}
\usepackage{setspace}
\doublespacing

\title{The iliac crest is an unreliable level landmark in transitional lumbosacral
       anatomy: a morphometric study of 802 screening CT examinations}

\author[1]{[First author]}
\author[1]{Gregory Schwing}
\author[1]{[co-authors]}
\affil[1]{[VERIFY affiliation]}

\date{}

\begin{document}
\maketitle

\begin{abstract}
\noindent\textbf{Introduction.} Surgeons locate the lumbar spine by palpating the iliac
crest, which is taught to lie at the L4--5 interspace. Lumbosacral transitional vertebrae
(LSTV) make level assignment ambiguous by altering the vertebral count. Whether the
palpable landmark remains reliable in those same patients has not been quantified.

\noindent\textbf{Methods.} 802 abdominal CT examinations from a colorectal cancer
screening cohort (393 women, 345 men; median age 59, range 50--89) were segmented for
every vertebra, rib, hip and femur. For each case we recorded the vertebral level reached
by the highest point of the iliac crest, and the number of rib-free vertebrae between the
lowest rib and the sacrum --- a count-free description of transitional status that does
not require naming a level.

\noindent\textbf{Results.} In spines with the typical five rib-free vertebrae, the crest
reached L4 in [X]\% of cases. In spines with four rib-free vertebrae it reached L4 in
only [X]\%, reaching L5 in [X]\% and L3 in [X]\%. [Add the association statistic once a
co-author has chosen the test.]

\noindent\textbf{Conclusions.} The iliac crest is a reliable indicator of the L4--5
interspace in typical anatomy and considerably less reliable in transitional anatomy. The
two sources of level error do not cancel: the landmark is least trustworthy in precisely
the patients in whom the count is already uncertain. Intraoperative imaging confirmation
should not be waived on the basis of a palpable crest when transitional anatomy is
suspected.
\end{abstract}

\section{Introduction}

% [CO-AUTHOR: expand. The pieces that need to be here:]
%  - LSTV prevalence and its range by definition (4-30%), and why the range is that wide
%  - wrong-level spine surgery: incidence, consequences, medicolegal weight
%  - the iliac crest / Tuffier's line as a taught landmark, for surgery and for neuraxial
%    anaesthesia, and the existing evidence that it is imprecise even in normal anatomy
%  - what has NOT been asked: whether landmark reliability and count ambiguity are
%    independent, or whether they fail together

Level identification in the lumbar spine rests on two independent things: a count, and a
palpable landmark. Both are known to be imperfect. What has not been established is
whether their errors are independent --- and if they are not, a surgeon relying on the
landmark to resolve an ambiguous count is relying on the wrong thing at the wrong moment.

\section{Materials and Methods}

\subsection{Cohort}

Records derive from the publicly available CTSpinoPelvic1K collection, itself built from
the TCIA CT~COLONOGRAPHY imaging archive with vertebral and pelvic segmentations unified
under a single label scheme. All examinations were performed for colorectal cancer
screening; no participant was imaged for a spinal indication. The lower age bound of 50
reflects the screening guideline rather than any selection applied here.

Ethical approval: the source collections are de-identified and publicly released;
[VERIFY institutional determination].

\subsection{Measurements}

All measurements are automated and deterministic; the code is public [repository URL].

\emph{Iliac crest level.} The most superior point of either iliac crest was located, and
the vertebral level whose segmentation spans that height was recorded.

\emph{Transitional status, without naming a level.} We counted the vertebrae lying
between the lowest rib-bearing vertebra and the sacrum. This is an interval count rather
than a level assignment and remains valid without deciding whether a given segment is a
sacralised L5 or a lumbarised S1 --- a distinction a spine-limited field of view cannot
make. Five is typical; four and six indicate transitional anatomy.

\emph{Quality control.} Every case passed geometric invariants: label geometry agreeing
with its CT in shape, affine and voxel spacing; no identifier outside the published
scheme; and left- and right-sided structures on the correct sides of the midline. Rib
numbering was verified by rib--vertebra incidence across 11{,}548 ribs, with 2 residual
offsets (0.017\%), both field-of-view truncations.

\subsection{Analysis}

% [CO-AUTHOR: choose and justify. Note this paper MAY test -- unlike the dataset article.]
% Proportions with Wilson score intervals; association between crest level and rib-free
% count by [chi-square / Fisher exact]; report effect size, not only a p-value.

Proportions are reported with Wilson score intervals, which remain within $[0,1]$ at the
small counts present in the transitional subgroups.

\section{Results}

\begin{table}[h]
\centering
\caption{Vertebral level reached by the iliac crest, by rib-free vertebral count.
Percentages with 95\% Wilson score intervals.}
\label{tab:crest}
\begin{tabular}{lccc}
\toprule
Rib-free count & $n$ & Crest at L4 & Crest elsewhere \\
\midrule
5 (typical) & [n] & [X]\% & L3 [X]\%, L5 [X]\% \\
4 & [n] & [X]\% & L3 [X]\%, L5 [X]\% \\
6 & [n] & [X]\% & L3 [X]\%, L5 [X]\% \\
\bottomrule
\end{tabular}
\end{table}

Cases in which the crest reached a thoracic level are excluded from
Table~
ef{tab:crest}: in those examinations the pelvis was barely within the field of
view, which is a coverage fact about the scan rather than an anatomical one about the
patient.

The direction of the shift differs between the two transitional groups. Where there are
four rib-free vertebrae the crest moves in both directions, appearing at L3 and at L5 in
roughly equal proportion; where there are six it moves predominantly caudally. Both
groups are small and neither proportion should be read as a precise estimate --- the
intervals in Table~
ef{tab:crest} are wide for that reason. What is robust is the
collapse in reliability itself, since the interval for the typical group does not
approach that of the four-vertebra group.

% [FIGURE: grouped bar chart with Wilson intervals, exported from the gallery]

\section{Discussion}

The clinical reading is a single sentence: \emph{the landmark is least reliable in exactly
the patients whose count is already uncertain.} In typical anatomy the crest identifies
L4 in roughly seven of eight patients; in the presence of four rib-free vertebrae it does
so in fewer than half, and it is as likely to point at L5 as at L3. Level-identification error has two
sources, and a surgeon who resolves an ambiguous count by palpating the crest is
compounding them rather than cross-checking one against the other.

% [CO-AUTHOR: argue with this. Points that need making, and points that need attacking:]
%  - this is a SCREENING cohort, not a surgical one. Selection differs. Say so.
%  - the crest level was measured on supine CT; palpation is done prone or lateral, and
%    soft tissue lies between finger and bone. The CT measurement is an upper bound on
%    landmark accuracy: palpation cannot be MORE accurate than the bone it is feeling for.
%  - transitional status here is by rib-free count. Radiographic LSTV classification
%    (Castellvi) is a different axis and is not interchangeable.
%  - what follows practically: when transitional anatomy is suspected, the crest should
%    not substitute for intraoperative imaging.

\subsection{Limitations}

The cohort is asymptomatic and undergoing screening, so these distributions are closer to
a reference range than a surgical series would be. Crest level was determined from supine
CT; palpation adds soft-tissue error on top of the bony variation measured here, so these
figures are an upper bound on landmark reliability rather than an estimate of it.
Transitional status is described by rib-free count rather than by Castellvi grade, and the
two are not interchangeable.

\section*{Data availability}
All measurements, the segmentations they derive from, and the code are publicly available
[DOI].

\section*{Conflicts of interest}
% [VERIFY]

\begin{thebibliography}{9}
% [ADD] Castellvi 1984; Bron 2007 (LSTV prevalence); Konin & Walz 2010 (AJNR, LSTV
% [ADD] classification and clinical relevance); wrong-level surgery incidence; Tuffier's
% [ADD] line accuracy studies; Vialle 2005; the CTSpinoPelvic1K dataset article
\end{thebibliography}

\end{document}
